\documentclass{article}
\usepackage{fontspec}
\usepackage{multicol}
\setlength{\columnsep}{1cm}
\setmainfont{Libertinus Serif}
\usepackage[a4paper, margin=1in]{geometry}
\usepackage{amsmath, amsthm, amssymb, mathtools, enumitem, hyperref}
\usepackage[portuguese]{babel}
\usepackage{csquotes}
\usepackage[backend=biber, style=numeric, sorting=nty]{biblatex}\addbibresource{referencias.bib}

\theoremstyle{plain}
\newtheorem{thm}{Teorema}[section]
\newtheorem{prop}[thm]{Proposição}
\newtheorem{lem}[thm]{Lema}
\newtheorem{cor}[thm]{Corolário}

\theoremstyle{definition}
\newtheorem{defn}{Definição}[section]
\newtheorem{exmp}[defn]{Exemplo}
\newtheorem{prblm}[defn]{Problema}

\theoremstyle{remark}
\newtheorem*{rem}{Nota}

\DeclareMathOperator{\spect}{spect}
\DeclareMathOperator{\ind}{ind}

\begin{document}
\begin{center}
    \textbf{\large UNIVERSIDADE FEDERAL DO RIO GRANDE DO SUL} \\
    \textbf{\large INSTITUTO DE MATEMÁTICA E ESTATÍSTICA} \\
    \vspace{3cm}
    \textbf{\Large RELATÓRIO DE PROJETO DE INICIAÇÃO CIENTÍFICA} \\
    \vspace{1cm}
    \textbf{\large TÉCNICAS COMBINATÓRIAS E ANALÍTICAS EM TEORIA DE GRAFOS} \\
    \vspace{3cm}

    \begin{tabular}{rl}
        \textbf{Acadêmico:} & LEONARDO KAWAI SUAREZ \\
        \textbf{Orientador:} & Prof. Dr. VILMAR TREVISAN \\
    \end{tabular}

    \vspace{12cm}
    \textbf{\large Porto Alegre, RS} \\
    \textbf{\large \today}
\end{center}

\newpage

\tableofcontents
\newpage

\section{Matriz de adjacência}



\section{Matriz de incidência}
\section{Matriz Laplaciana}
\section{Matriz Laplaciana sem sinal}
\section{Algoritmos de Localização de Autovalores}

\newpage

\printbibliography[title={Referências}]

\end{document}